\documentclass[10pt]{article}
\usepackage[a4paper,margin=25mm]{geometry}
\usepackage{amsmath,amssymb,amsthm,mathrsfs,longtable,hyperref}
\setlength{\emergencystretch}{3em}
\newtheorem{theorem}{Theorem}[section]
\newtheorem{lemma}[theorem]{Lemma}
\newtheorem{proposition}[theorem]{Proposition}
\newcommand{\PSH}{\operatorname{PSH}}
\newcommand{\FS}{\operatorname{FS}}
\newcommand{\MA}{\operatorname{MA}}
\newcommand{\vol}{\operatorname{vol}}
\newcommand{\EE}{\mathcal E^1_\theta}
\newcommand{\OO}{\mathcal O}
\title{Original-carrier finite-energy equations in big rational classes\\
\large A reconstruction and audit of the supplied moving-system programme}
\author{Unverified research draft}
\date{6 October 2026}
\begin{document}
\maketitle
\begin{abstract}
We reconstruct the retained first domain of original P16.5: every actual
big rational bundle on every smooth integral projective variety over every
algebraically closed trivially valued characteristic-zero field. Only G0 in
its actual ample domain and G1 for relative generation and one FS net are
caller grants. We construct the original all-point cone, its genuine minimal
reference and entire finite-energy space. Actual section-volume derivatives,
finite intersection limits, quotient traces and bounded physical floors
construct the genuine mixed operator. Finite truncation gives unbounded
energy-complete approximation independence without output convergence in
the admissibility condition. Atomic dual density precedes attainment and
proves weak target-integral continuity. Compactness and both derivative
signs give actual equations; whole planes and divisorial probes give
every-solution uniqueness. Fixed ramps give strict Borel comparison.
The additional zero-volume range proved is exactly the numerically trivial
rational-bundle range. A stronger fixed original unbounded-test assertion
at every physical auxiliary output is explicitly not claimed or used.
This is an unverified proof draft, not a correctness certification.
\end{abstract}

\section{Task, scope and authority}
The unchanged first target of P16.5 is every actual big rational line bundle
on every smooth integral projective variety over every algebraically closed
trivially valued characteristic-zero field. Throughout, $X$ has dimension
$n\geq1$, $K=X^{\rm an}$, $L\in\operatorname{Pic}(X)\otimes\mathbb Q$
is actual and big, $\theta=c_1(L)$, and $V=\vol_X(L)>0$.
The generic trivial valuation is $v_0$. All supported semivaluations are
retained. There is no countability, invariant, skeletal, support or
section-ring finite-generation hypothesis. All section degrees are cleared
before sections, ideals and tensor products occur.
The coordinate set is the rationally scaled divisorial valuations together
with $v_0$. A finite ``divisorial'' model measure below is allowed to have
an atom at $v_0$; that coordinate is always retained. This convention is
necessary already for the trivial model and its mass-$V$ point measure.

Only the complete supplied G0 and G1 bodies are mathematical grants.
G0 is used on actual ample rational bundles, including their genuine cone,
finite intersections, compactness, all-point divisorial determination and
finite ample locality. Its equation is not used to construct a big-class
operator or a target approximation. G1 is used only on its resolved smooth
equivariant actual models for its relative multiplier-ideal generation and
the one decreasing FS net. All other supplied papers and ledgers are proof
data. We neither open the historical paths in those ledgers nor infer a
theorem from a historical label or hash.

Ordinary algebraic background is characteristic-zero projective equivariant
resolution and principalization, Kodaira decomposition, numerical openness
of ampleness, high-power generation, proper normal birational section
descent, section-volume asymptotics and Fujita approximation. BFJ's
positive products and volume differential are used on smooth projective
algebraic determinations in their published characteristic-zero domain.
Their complete pertinent units were independently retrieved. These are
algebraic inputs, not nonample non-Archimedean equation grants.

The theorem below concerns the entire retained first big domain.
The actual nonample foundations are reconstructed in Sections 2--10,
rather than assumed as S--M. The numerically trivial zero-volume
range is stated separately. Nonzero nonbig pseudoeffective classes
are not asserted to have the same raw top equation or uniqueness.
No other P16 axis is claimed solved. The exact original task and full
family remain in the frozen input.

\section{The original cone and its first envelopes}
An actual full generically finite FS scalar on $B$ is
\[
 h=m^{-1}\max_j\{\log\|s_j\|_{\rm triv}+\lambda_j\},
 \qquad 0\ne s_j\in H^0(X,mB),\quad\lambda_j\in\mathbb R.
 \tag{2.1}
\]
Common zeros are permitted. Define $\PSH(\theta)$ to be all-point
decreasing endpoints of one such net on actual $B_i$ with
$c_1(B_i)\to\theta$, whose generic endpoint is finite. This definition
does not assume closure of a twice-closed family. Weak topology means
pointwise convergence at rational divisorial valuations and $v_0$.
Every scalar in (2.1) has its supremum at $v_0$: a canonical section norm
is at most one everywhere and equals one at $v_0$.

\begin{lemma}[Literal ample embedding]\label{embed}
Choose an actual ample $\Omega$ with $\Omega-L$ ample. Every genuine
$\theta$ scalar is the identical all-point function in G0's ample cone
for $\Omega$. Its divisorial values are finite and determine all-point
order. For each fixed divisorial $z$ there is a finite $C_z$ such that
$-C_z\leq u(z)\leq0$ for every normalized $u\in\PSH(\theta)$.
\end{lemma}
\begin{proof}
On a tail $\Omega-B_i$ is ample. Tensor with a weight-zero generating
system of a cleared power of that actual difference. The scalar is
unchanged: some generator is a unit at the valuation-ring center of every
analytic point. Append generators of a power of $\Omega$ at weights
$-mN$. The product-directed net has scalar $\max(h_i,-N)$, decreases,
is on the actual fixed $\Omega$, and has endpoint $u$ everywhere.
Thus G0 applies. The coordinate bounds follow by continuity of the finite
coordinate on its compact normalized ample cone. No energy comparison
between $L$ and $\Omega$ is used.
\end{proof}

Finite maxima have one defining net: place the incoming section systems
on $L+\varepsilon A$ using actual ample differences and take their union.
Valid indices are directed by a smaller rational $\varepsilon>0$ and
later incoming indices. For real $0<r<1$, put incoming scalars below $C$
and take rational $r_j\uparrow r$, $s_j\uparrow1-r$. The scalars
$r_j(h_i-C)+s_j(g_l-C)+C$ decrease on the product index, are actual FS
after clearing powers, and have the desired convex endpoint, including
at poles. Constants and integer tensoring are immediate.

Fix an actual Kodaira identity
\[
 L=aA+E_0,\qquad a>0,\quad A\text{ ample},\quad E_0\geq0,
 \quad W=\log\|s_{E_0}\|\leq0,
 \quad Z=E_0^{\rm an},\quad U=K\setminus Z.
 \tag{2.2}
\]
The scalar $W$ is an actual FS scalar on $L$ by tensoring its section
with a generated power of $aA$. It is finite continuous on $U$.

\begin{lemma}[Section minorants before closure]\label{density}
For bounded continuous $f$ and genuine $u\leq f$,
$u\leq Q_L(f):=\sup\{h\in\FS(L):h\leq f\}$ on $U$.
\end{lemma}
\begin{proof}
Let $h_i\downarrow u$ be its literal incoming net on $B_i$.
Fix rational $0<t<1$ and $\varepsilon>0$ first. The actual difference
$L-(1-t)B_i-t(L-aA)$ is eventually ample. The nested compact sets
$\{h_i\geq f+\varepsilon\}$ have empty intersection; hence a tail
has $h_i<f+\varepsilon$ everywhere, also at supported points.
Tensoring into the displayed difference gives the actual $L$ minorant
$(1-t)(h_i-\varepsilon)+t(W-\|f\|_\infty)\leq f$.
At $x\in U$ take the incoming endpoint, then $t\downarrow0$, then
$\varepsilon\downarrow0$. If $u(x)=-\infty$ the conclusion is immediate.
At a supported pole all section inequalities remain literal extended
inequalities; no infinite values are subtracted.
\end{proof}

For a rational continuous model obstacle $f$, G1 supplies one genuine
endpoint $\Psi$ with $Q_L(f)\leq\Psi\leq f$. Lemma~\ref{density}
puts every genuine competitor below $\Psi$ on $U$, in particular at all
divisorial points. Lemma~\ref{embed} then proves that inequality at every
original point. Thus $\Psi=P_\theta(f)$, the raw all-point envelope.
In particular
\[
 R=P_\theta(0)\in\PSH(\theta),\qquad R(v_0)=0,
 \qquad u\leq R\leq0\quad(u(v_0)=0).
 \tag{2.3}
\]
The same construction works for every actual big rational $B$.
A fixed late G1 stage, normalized to $\psi_B\leq0$ on a bundle close
to $B$, gives a finite constant $c_B$ with $R_B\leq\psi_B+c_B$.
Separately normalized stages are not claimed to form a decreasing net.
For every bounded obstacle the constant envelope law also gives
$R+\min f\leq P_\theta(f)\leq R+\max f$. Thus these model envelopes
have bounded relative singularity difference; after energy is defined
they are actual finite-energy functions by its monotonicity and constant law.

\section{Single-net closure, envelopes and compactness}
\begin{lemma}[Old pole comparison]\label{pole}
If actual FS scalars satisfy $\psi\leq0$ and
$F\geq(1-\delta)\psi-k$, $0<\delta<1$, and usc $h\leq\psi+C$,
then $h-F$ off the poles of $F$, assigned $-\infty$ at those poles,
is usc. Decreasing such $h$ give a decreasing comparison net. If its
infimum is negative everywhere, an entire tail is negative everywhere.
\end{lemma}
\begin{proof}
Off the poles its upper bound is $\delta\psi+C+k$. At a pole of $F$
the inequality forces $\psi=h=-\infty$; extended continuity of $\psi$
makes that bound tend to $-\infty$. The assignment is a comparison
extension, not the value of $(-\infty)-(-\infty)$. Nested compact
nonnegative superlevels and directedness prove the last assertion.
\end{proof}

\begin{theorem}[Actual decreasing closure]\label{closure}
A decreasing net $u_i\in\PSH(\theta)$ with finite generic infimum
has one all-point decreasing FS defining net with classes tending to
$\theta$.
\end{theorem}
\begin{proof}
Work on a tail with $u_i\leq M$. For rational $\eta>0$ put
$T_\eta=L+\eta A$, choose $0<\epsilon<\eta/4$, $B=L+\epsilon A$,
and a fixed normalized G1 reference stage $\psi$ on $S$ so close to
$B$ that $T_\eta-S$ is ample. Choose small rational $\delta>0$ with
$T_\eta-(1-\delta)S$ ample. A special stage is an FS scalar on
$T_\eta$ with complete such data and
\[
 F\geq u+\kappa,\qquad F\geq(1-\delta)\psi-k,
 \qquad\kappa>0.
 \tag{3.1}
\]
It exists by taking the maximum of a shifted late defining majorant
of some $u_i$ and the floor, tensoring both into $T_\eta$.
For every fixed old $B_l,\psi_l$, the $u_i$ and late defining stages
$h$ of a fixed $u_i$ embed into $B_l$. Their generic values are at
most $M+1$, so first envelope maximality gives
$h\leq R_{B_l}+M+1\leq\psi_l+c_l+M+1$. This also bounds all $u_i$.

Given finitely many old stages $F_l$, choose the new $\eta$ arbitrarily
small below their $\epsilon_l$. Choose the new $S$ close enough to
$L+\epsilon A$ that every $B_l-S$ is ample. Then
$\psi\leq R_{B_l}\leq\psi_l+c_l$. Take
$0<\delta<\min_l\delta_l$. Where $F_l$ is finite, the new floor
$b=(1-\delta)\psi-k$ obeys
$b-F_l\leq(\delta_l-\delta)\psi_l+(1-\delta)c_l+k_l-k$.
One large $k$ makes it at most zero. At an old pole $b=-\infty$.
Choose $0<\kappa<\min_l\kappa_l/4$. Lemma~\ref{pole} first applied
to $u_i+2\kappa$ selects one late $i$ below every old strict guard.
Applied again to its literal defining tail, with the same old bounds,
it selects $h$ fitting the new target and with $h+\kappa\leq F_l$.
Thus $\max(h+\kappa,b)$ is a new special stage below all old ones.
Order all special stages by scalar and $\eta$ decrease. They are directed.

For $x$ and finite $t>u(x)$ choose a late $u_i(x)<t$ with margin,
then its target gap before a late defining $h$ fitting that gap, then
a small guard and large floor constant so both branches are below $t$
at $x$. This also works for every finite $t$ if $u(x)=-\infty$.
Subsequent directed refinements impose every old constraint. All stages
are above $u$, so their all-point infimum is exactly $u$. Their classes
are $\theta+\eta[A]\to\theta$ and their generic endpoint is finite.
\end{proof}

Every upper bounded usc $f$ on compact Hausdorff $K$ is the pointwise
infimum of continuous majorants: use a cutoff inside an open strict
sublevel at the tested point and a global upper bound outside. Uniform
rational model density and finite minima make these model majorants
directed and decreasing. If $f$ has a genuine minorant of finite generic
value, their first envelopes decrease with that same generic lower bound.
Theorem~\ref{closure} makes their endpoint genuine and exactly $P_\theta(f)$.
For an upper bounded genuine family, its raw supremum $S$ has usc
regularization $S^*$, and $P_\theta(S^*)\geq S$ is usc, hence dominates
$S^*$. The reverse obstacle inequality proves equality all-point.

\begin{theorem}[Original normalized compactness]\label{compact}
The genuine normalized $\theta$ cone is compact in its divisorial topology.
\end{theorem}
\begin{proof}
Embed into G0's compact normalized actual ample cone. For a net in the
big cone with ample weak limit $a$, form its regularized tail suprema
$S_i^*=(\sup_{j\geq i}u_j)^*$. They are genuine in the big cone by
the preceding envelope proof, decrease and have generic value zero.
Their endpoint $b$ is genuine by Theorem~\ref{closure}.
The independently read primary ample separator theorem (BJ 5.6--5.7)
gives $S_i^*(z)=\sup_{j\geq i}u_j(z)$ at each divisorial $z$.
Thus $b(z)=a(z)$, and G0's all-point determination gives $b=a$.
The big cone is a closed subset of the compact ample cone.
Rational scalings suffice: on each divisorial ray the FS stages are
convex maxima of affine functions; finite ample coordinate bounds give
uniform secant Lipschitz bounds on smaller intervals. Rational values
therefore determine real-scaled values and their convergence.
\end{proof}

\section{Full-section scalar energy and finite model derivatives}
On a smooth equivariant modification $Y\to X\times\mathbb P^1$,
product away from zero, let $F=\sum_E b_EE$ be the full fiber, and
$D$ vertical rational Cartier. Put $f=f_D$, $f(v_E)=d_E/b_E$,
$v_E=b_E^{-1}\operatorname{ord}_E|_{k(X)}$.
If $aF\leq D\leq bF$, the full actual relative lattice is
\[
 \Lambda_m=\bigoplus_jt^{-w_{mj}}k[t]s_{mj},\quad
 ma\leq w_{mj}\leq mb,\quad (s_{mj})\text{ a basis of }H^0(X,mL).
 \tag{4.1}
\]
Divisor comparison, $\rho_*\OO_Y=\OO_{X\times\mathbb P^1}$ and
projection formula give its bounding lattices. Equivariance grades it;
an adapted basis of the increasing homogeneous flag gives (4.1).
Gluing to the product lattice at infinity gives
\[
 h^0(Y,m(\rho_X^*L+D+cF))
       =\sum_j\max(w_{mj}+mc+1,0).
 \tag{4.2}
\]
For $c>-a$, full-section volume asymptotics imply
\[
 \vol_Y(\rho_X^*L+D+cF)/(n+1)=cV+\mathscr E(f),\qquad
 \mathscr E(f)=\lim_m\frac{n!}{m^{n+1}}\sum_j w_{mj}.
 \tag{4.3}
\]
These are actual bundles and complete systems. The intercept is invariant
under refinement, monotone, $V$-translation equivariant and $V$-Lipschitz
for the original sup norm. Expanding algebraic volume-root concavity in
(4.3) at $c\to\infty$ proves its ordinary concavity.

BFJ differentiation on the actual big $(n+1)$-fold gives
\[
 \mu_D=\sum_E b_E\langle(\rho_X^*L+D+cF)^n\rangle E\,\delta_{v_E},
 \quad\mu_D(K)=V,\quad D\mathscr E_f(h)=\int h\,d\mu_D.
 \tag{4.4}
\]
Every coefficient is nonnegative, since increasing the effective $E$
cannot decrease volume. Differentiating (4.3) in $D+sE$ on a common
eventual half-line proves independence of large $c$ coefficient by
coefficient. Differentiation in $F$ gives the mass. On a refinement an
effective exceptional prime $G$ has $p_*\OO(rG)=\OO$ for $r\geq0$;
adding it does not change sections, so its derivative coefficient is zero.
The remaining coefficients agree by pullback and projection.

For continuous $f$ and $s\ne0$ in the full $R_m=H^0(X,mL)$ set
\[
 W_m(f;s)=\inf_K(mf-\log\|s\|).
 \tag{4.5}
\]
At a section zero its bracket is $+\infty$. Its weights lie in
$[m\min f,mf(v_0)]$. Superlevel subspaces form a finite closed-endpoint
flag: an intersection of a decreasing family of subspaces of a fixed
finite-dimensional space stabilizes. Its trace is basis independent.
An adapted basis with floored weights has scaled trace $e_m(f)$.
Uniform flag shifts give
$|e_m(f)-e_m(g)|\leq(n!\dim R_m/m^n)(\|f-g\|+1/m)$.
For a model obstacle its flag is exactly (4.1), by the all-point section
extension calculation in G1. Model density and rank asymptotics give a
unique limit $\mathscr E(f)$ for every continuous original obstacle.

The section inequality in (4.5) is equivalent to its weighted section
minorant lying below $P_\theta(f)$. Therefore any continuous
$g\geq P_\theta(f)$ has $\mathscr E(g)\geq\mathscr E(f)$.
Define on the entire actual cone
\[
 E(u)=\inf_{f\in C(K),\,f\geq u}\mathscr E(f),\qquad
 \EE=\{u\in\PSH(\theta):E(u)>-\infty\}.
 \tag{4.6}
\]
Then $E(P_\theta(f))=\mathscr E(f)$, $E(R)=0$ and
$E(u+c)=E(u)+cV$. The reference is the genuine actual $R$, not an
arbitrary lift or a numerical replacement for the bundle.

For $u_i\downarrow u$ and fixed continuous $f\geq u$, the compact
sets $\{u_i\geq f+\varepsilon\}$ eventually vanish. Hence
$\inf_i E(u_i)\leq\mathscr E(f)+\varepsilon V$, proving decreasing
continuity, including at $-\infty$. Relative floors $\max(u,R-N)$
are genuine and their energies lie in $[-NV,0]$ for normalized $u$.
Thus (4.6) really uses the entire cone, not only functions with bounded
relative difference. Nonpositive rational model majorants normalized at
$v_0$ are obtained by taking the minimum of a model majorant and zero;
their envelopes decrease all-point with decreasing energies.

Finite concavity gives global planes
\[
 E(u)\leq A_D(u):=\mathscr E(f_D)+\int(u-f_D)d\mu_D,
 \qquad E(u)=\inf_D A_D(u).
 \tag{4.7}
\]
Pass continuous majorants on the fixed finite divisorial support for the
first inequality. For equality use model majorants of $u$, whose plane
correction is nonpositive, and then their outer infimum. Each plane is
weakly continuous since all its divisorial coordinates are finite.
Consequently $E$ is weakly usc and concave, and normalized sets
$E\geq-B$ are compact sets of actual finite-energy functions.

\section{Complete moving systems and the retained trace error}
For any big actual $N$ on a smooth projective $d$-fold, resolve the full
$|mN|$ and write $p_m^*(mN)=M_m+G_m$, $A_m=M_m/m$ generated,
$G_m\geq0$. Multiplication gives $A_{m\ell}\geq A_m$ on common
determinations. Cofinality in every positive power follows directly:
for finitely many nef competitors $\beta_i\leq N$ choose actual ample
$Q_0$ with $N-Q_0$ big. Then $Q_i=(1-\epsilon)\beta_i+\epsilon Q_0$
is ample and $N-Q_i$ big. In one cleared degree all $qQ_i$ generate
and all $q(N-Q_i)$ have sections with divisors $R_i$. Their injections
into the whole $|qN|$ give $G_q\leq R_i$ and $A_q\geq Q_i$.
Products of nef classes with a psef difference have the correct order.
Let $\epsilon\downarrow0$ on this fixed determination. BFJ's directed
compactness converts the equality of upper bounds into
\[
 A_m^p\longrightarrow\langle N^p\rangle\quad(1\leq p\leq d).
 \tag{5.1}
\]
Real competitors follow by interior approximation. A fixed signed
Cartier test is first written as the difference of two fixed ample tests;
pair (5.1) separately and subtract. No positive-product multilinearity
has been inferred from nef polynomial formulas.

For fixed finitely many model obstacles $f_D\leq0$ on $X$, the complete
generic system gives the actual identity $p_m^*L=A_m+T_m$, with
$mA_m$ generated and $mT_m$ effective. A sufficiently large fixed time
twist $c$ has the entire absolute basis
$t^{-(w_{mj}+mc)+r}s_{mj}$, $0\leq r\leq w_{mj}+mc$.
Its infinity endpoints recover every original $s_{mj}$; all nonzero
finite fibers also recover the full generic system. Remove the horizontal
factor $s_{mT_m}$ before resolving the residual ideal. Its residual base
divisor is vertical, and its moving bundle is
\[
 C_{m,D}=\rho_X^*A_m+D+cF-G_{m,D}^{\rm vert}/m.
 \tag{5.2}
\]
It is globally generated and nef. The actual physical section identity is
\[
 \phi_{m,D}=q_{T_m}+\psi_{m,D},\qquad\phi_{m,0}=q_{T_m},
 \quad\psi_{m,D}\leq0.
 \tag{5.3}
\]
These statements retain all horizontal fixed factors and every $b_E$.
They do not give original finite energy to $\phi_{m,D}$.
Applying (5.1) on $X$ and the total model gives
\[
 V_m=A_m^n\to V,\quad
 e_m(\psi_{m,D})=C_{m,D}^{n+1}/(n+1)-cV_m\to E(P_\theta(f_D)),
 \quad p_{m*}\MA_{A_m,\rm raw}(\psi_{m,D})\to\mu_D.
 \tag{5.4}
\]
The last limit pairs (5.1) with each fixed vertical Cartier test, then
uses model density and the converging raw masses. Finite nef model
measures have coefficient $b_E(C_{m,D}|_E)^n$ by G0's finite ample
perturbation, with the model fixed through that polynomial limit.
The divisibility-increasing full section stages have the envelope as
their supremum on $U$ by Lemma~\ref{density}. A cleared Kodaira section
times a generating ample system makes a cofinal tail finite on $U$.
The decreasing G1 multiplier-ideal stages are finite there too: their
base ideals contain the corresponding full-system base ideals, and their
additional fixed model weight is bounded. Thus the same endpoint on $U$
is both a supremum and an infimum of finite continuous functions; it is
finite continuous. Compact positive-gap superlevels now give the net
Dini argument. It proves compact-uniform convergence of $q_{T_m}$ to $R$
and of the quotient to $P_\theta(f_D)-R$.
Their generic values tend to $P_\theta(f_D)(v_0)$ by the same section
minorants at $v_0$. This is zero for the normalized envelopes used in
the subsequent zero-generic-error limits. A general nonpositive obstacle
need not have normalized envelope; its finite stages are shifted by the
same actual constant when normalization is required.

For fixed $m$, a bounded nonpositive relatively generated scalar $\psi$
on $A_m$, and a fixed original continuous $f\geq q_{T_m}+\psi$,
multiply by the actual $s_{\ell T_m}$ to inject the full auxiliary
$H^0(\ell A_m)$ into $H^0(\ell L)$. Injected weights are at least the
auxiliary metric weights; a supported zero contributes $+\infty$ to
the original infimum. All other original weights are at least
$\ell\min f$. Extend an adapted subspace basis to the full filtered
basis and take $\ell\to\infty$ at fixed $m,f,\psi$. Rounding errors
are bounded by the ranks. This proves
\[
 \mathscr E(f)\geq e_m(\psi)+(\min_K f)(V-V_m).
 \tag{5.5}
\]
The last, generally negative, term is retained until $m\to\infty$ at
this fixed finite $f$, before the final outer infimum.

\section{Finite Hodge estimates and genuine contraction}
On a smooth SNC total model with $F=\sum_i b_iE_i$, relatively nef
slots $\Gamma=N_1\cdots N_{n-1}$, and $D=\sum_i b_i a_iE_i$,
\[
 -D^2\Gamma=\sum_{i<j}b_ib_j(E_iE_j\Gamma)(a_i-a_j)^2\geq0.
 \tag{6.1}
\]
This follows from the zero row identities
$\sum_j b_j(E_iE_j\Gamma)=0$ and positivity on each SNC stratum.
Resolve $|f_D|$ by principalizing
$\OO(-q(CF+D))+\OO(-q(CF-D))$ after clearing denominators.
Its normalized coefficients are exactly $|a_i|$ on the refined model.
The inequality $||a_i|-|a_j||\leq|a_i-a_j|$ and projection formula
prove contraction of the quadratic form under absolute value; the same
proof works for rational piecewise linear one-Lipschitz functions.

On an actual generated nef $A$, put $T_\xi=(A,\xi)$,
$e_A(\xi)=T_\xi^{n+1}/(n+1)$, $J_A(\xi)=V_A\xi(v_0)-e_A(\xi)$.
For $S=(0,\xi-\eta)$ the correct finite tangent formula is
\[
 T_\eta^{n+1}-T_\xi^{n+1}+(n+1)ST_\eta^n
 =-\sum_{j=0}^{n-1}(n-j)S^2T_\xi^jT_\eta^{n-1-j}.
 \tag{6.2}
\]
Expanding two telescoping sums gives the multiplicity $n-j$.
Let $d_A(\xi,\eta)$ be the maximum of these $n$ quadratic forms.
Then $J_A(\xi)$ is their weighted sum relative to zero divided by
$n+1$, and hence is comparable to $d_A(\xi,0)$ with dimension constants.
The midpoint defect $\delta_A$ satisfies
\[
 \frac{n}{4^{n+1}}d_A(\xi,\eta)\leq\delta_A(\xi,\eta)
 \leq\frac n8 d_A(\xi,\eta).
 \tag{6.3}
\]
For $p(t)=e_A(\eta+t(\xi-\eta))$, $-p''$ is $n$ times the positive
Bernstein combination of those forms. Integrate against
$\min(t,1-t)/2$. Its integral is $1/8$; on $[1/4,3/4]$ the kernel
and each Bernstein monomial give the stated lower coefficient.

The finite slot-change calculation uses only (6.1), Cauchy--Schwarz and
(6.2). For completeness set $\tau=(\xi+\eta)/2$, $f=\xi-\eta$,
$a=d_A(\xi,\eta)$, $b=\max(d_A(\xi,\rho),d_A(\eta,\rho))$,
$b_j=\|f\|^2_{T_\rho^jT_\tau^{n-1-j}}$.
The average expansion gives $b_0\asymp_n a$, while the seminorm
triangle inequality gives $b_{n-1}\leq4b$. Replacing one slot gives
$b_{j+1}-b_j=-S^2(0,\rho-\tau)T_\rho^jT_\tau^{n-2-j}$.
Write one $S=T_\xi-T_\eta$ and apply Cauchy--Schwarz separately in
the two nef forms. Since each $T_\xi,T_\eta$ is bounded by $2T_\tau$
in positive quadratic forms, this is at most $C_n\sqrt{b_jb}$.
Here quasi-convexity of $d_A$ follows from the convex tangent remainder
in (6.2). If $a\leq b$, induction gives
$b_j\leq C_n a^{2^{-j}}b^{1-2^{-j}}$; if $a>b$, use $b_{n-1}\leq4b$.
Applying this to differences through a third potential gives the
dimension-only quasi-triangle inequality. Average a finite desired list
of slots; its expanded power dominates each mixed positive quadratic
form. With $\beta=2^{1-n}$ this proves
\[
 \|\xi-\eta\|^2_{T_{\rho_1}\cdots T_{\rho_{n-1}}}
 \leq C_n d_A(\xi,\eta)^\beta
       \max(J_A(\xi),J_A(\eta),J_A(\rho_l))^{1-\beta}.
 \tag{6.4}
\]
This reconstructs the primary finite calculation with the corrected
tangent orientation. It imports no unbounded ample estimate.

Resolving $g=|\xi-\eta|$, telescope from the trivial measure
$V_A\delta_{v_0}$ to $\MA_A(\zeta)$, starting at $j=0$.
Apply Cauchy--Schwarz, absolute contraction and (6.4) in each term.
With $\alpha=2^{-n}$ and $B_* =\max(J_A(\xi),J_A(\eta),J_A(\zeta))$,
\[
 \int|\xi-\eta|d\MA_A(\zeta)
 \leq V_A|\xi(v_0)-\eta(v_0)|
   +C_n\delta_A(\xi,\eta)^\alpha J_A(\zeta)^{1/2}B_*^{1/2-\alpha}.
 \tag{6.5}
\]
All multiplicities and raw volumes are present. For $n=1$ omit the
last exponent-zero factor. Zero-volume cases follow from positivity and
mass; no division by $V_A$ occurs.

For $\psi\leq0$ set
$a_k=\int(-\psi)\MA_A(\psi^{[k]},0^{[n-k]})\geq0$.
The finite energy identity gives $-e_A(\psi)=(n+1)^{-1}\sum a_k$,
and $a_0=-V_A\psi(v_0)$. Expanding the finite nef polynomial gives
\[
 -e_A(r\psi)\leq r a_0+2^{n+1}r^2(-e_A(\psi)),\quad0\leq r\leq1.
 \tag{6.6}
\]
The terms with $k=0$ sum to at most $ra_0$; every other term has
another factor $r$, with total binomial coefficient bounded by $2^{n+1}$.
For $u=P_\theta(f_D)$ use the actual physical minorant
$q_{T_m}+r\psi_{m,D}\leq(1-r)R+ru$ in (5.5). Hold each continuous
majorant fixed, then take $m\to\infty$ in (5.4), retaining the vanishing
generic term and the missing-volume term, then take the outer infimum.
This proves the genuine contraction
\[
 B((1-r)R+ru)\leq2^{n+1}r^2B(u),\qquad B(u)=-E(u).
 \tag{6.7}
\]
Canonical model envelopes and decreasing energy continuity prove it for
every normalized actual $u\in\EE$. Rational $r$ extends to real $r$
by concavity and continuity along its genuine finite-energy segment.

For every normalized genuine $q$, the truncations
$q_j=\max(q,R-16^j)$ have cost at most $V16^j$.
Put $a_j=(1-8^{-j})R+8^{-j}q_j$ and
$p_m=\sum_{j\leq m}2^{-j}a_j+2^{-m}R$.
Contraction and concavity bound its cost by $2^{n+1}V/7$.
It decreases genuinely with generic value zero, so closure and decreasing
energy continuity give $p\in\EE$. Where $q=-\infty$ and $R$ is finite,
$R-p_m=m$. Thus this one actual E1 pole covers the pole of $q$ outside
the reference poles. In particular applying it to $W$ from (2.2)
constructs the required genuine finite-energy $Z$-pole.

\section{The original outer derivative and signed big margin}
For $f_D\leq0$, $\mathscr E(f_D)\geq-B$ and $c>B/V$, put
$g_D(c)=\vol(\rho_X^*L+D+cF)/(n+1)$.
Formula (4.2) makes $g_D$ convex, $V$-Lipschitz, zero at $c=0$ and
eventually $cV+\mathscr E(f_D)$, with $g_D(c)\geq cV-B>0$.
Thus the deficient class is genuinely big. Its positive derivative
measure $\mu_{D,c}$ has mass between $V-B/c$ and $V$.
There is coefficientwise domination
\[
 0\leq\mu_{D,c}\leq\mu_D,
 \quad\|\mu_D-\mu_{D,c}\|_{\rm TV}\leq B/c.
 \tag{7.1}
\]
Indeed the actual full-lattice trace
$\mathscr E(f_D)+cV-g_D(c)
=\lim n!m^{-n-1}\sum_j\min(w_{mj}+mc,0)$ is monotone under
adding any effective vertical divisor: inclusion of the full filtration
raises each ordered weight. Differentiate this monotonicity in each
effective $E$. The difference coefficient is nonnegative. Convexity
in $c$ gives its stated mass bound.

Fix one rational original model test $h=f_H$ on a fixed determination.
Choose ample $T$ there, $\gamma=\rho_X^*L+cF$, and
$\kappa=(n+1)(cV-B)>0$, $M=(n+1)\langle\gamma^n\rangle T$.
For every refinement carrying $D$, volume loss along $\alpha-rT$
while big is at most $rM$, since $\alpha=\rho_X^*L+D+cF\leq\gamma$.
Boundary continuity excludes a first exit before
$r_0=\kappa/[2(M+1)]$. Choose $a$ with $aT\pm H$ nef and put
$\lambda=a/r_0$. Then $\lambda\alpha\pm H$ are psef, so for
$|s|\lambda<1/2$ positive-product order sandwiches
$\langle(\alpha+sH)^n\rangle$ between
$(1\pm\lambda|s|)^n\langle\alpha^n\rangle$.
Pair separately with the two fixed nef terms $aT+H$ and $aT$.
This bounds the derivative variation uniformly over all $D$ and refinements
and gives a fixed $K_c$ quadratic remainder for $g_{D+tH}(c)$.

For the full energy use the nonpositive obstacle
$f_D+sh-|s|H_0$, $H_0=\|h\|_\infty$, and parameter
$c_s=c+|s|H_0$. Its cost is at most $B_s=B+2V|s|H_0$.
Impose, before differentiating either sign,
\[
 |t|H_0<(c-B/V)/2,\qquad |t|<1.
 \tag{7.2}
\]
Then $c_s>B_s/V$ throughout the path, whose actual deficient class is
exactly $\alpha+sH$. The rough bound $B+2VH_0$ by itself does not
prove this when $c$ is just above $B/V$. Integrate the positive derivative
tail (7.1) and restore the subtracted scalar. The result is
\[
 |\mathscr E(f_D+th)-\mathscr E(f_D)-t\mu_D(h)|
 \leq K_ct^2/2+|t|H_0(2B+2VH_0)/c.
 \tag{7.3}
\]
All order estimates use fixed nef tests. Both signed perturbations have
the verified margin (7.2).

For $u\in\EE$ use its canonical nonpositive model majorants $f_i$.
Their positive mass-$V$ measures have weak cluster points on original
compact $K$. At fixed $t$, their perturbed genuine envelopes decrease
exactly to $P_\theta(u+th)$, with decreasing energies. Pass the whole
increment (7.3) along a cluster, then divide with each sign, let $t\to0$
at fixed $c$, and finally $c\to\infty$. It gives a unique cluster $M(u)$
and
\[
 \left.\frac d{dt}\right|_0E(P_\theta(u+th))=\int h\,dM(u).
 \tag{7.4}
\]
Uniform model approximation changes each divided increment and integral
by at most $V$ times its sup error. Thus (7.4) holds for every original
continuous $h$. This is a target-independent outer-energy map, positive
of raw mass $V$, on actual functions on $K$.
The envelope and energy constant laws show directly from (7.4) that
$M(u+c)=M(u)$ for every real $c$; no target equation is used for this.
For ordered bounded-cost endpoints $a\geq b$, compare their perturbed
energy increments with (7.3); for fixed model $h$,
\[
 |M(a)(h)-M(b)(h)|\leq (E(a)-E(b))/t+K_ct+C_{h,B}/c.
 \tag{7.5}
\]
This proves decreasing weak continuity. At $u=P_\theta(f_D)$ apply
it to model majorants and their minima with $f_D$: the latter have
exactly the same envelope and energy as $f_D$. First (7.5) makes their
derivatives equal to $\mu_D$, then compares them to the canonical
majorants. Thus $M(P_\theta(f_D))=\mu_D$.

For an actual effective Cartier $S$, the vertical divisor moment of
$\mu_{D,c}$ is bounded by
$\langle\alpha_D(c)^n\rangle\rho_X^*S$; its effective horizontal
part has nonnegative volume derivative. A fixed nef upper test gives
one $C_{S,c}$ independent of $D$. For the original continuous caps
$b_N=\min(v(S),N)$, (7.1) implies
\[
 \int b_N\,dM(u)\leq C_{S,c}+NB/c.
 \tag{7.6}
\]
Divide by $N$, take $N\to\infty$, then $c\to\infty$.
Hence $M(u)(S^{\rm an})=0$, in particular $M(u)(Z)=0$.
This argument gives algebraic nullness and uniform algebraic collars.
It does not yet identify all unbounded genuine weighted products.

\section{Intrinsic targets and what pre-attainment estimates supply}
For a positive original Radon measure $\nu$ of raw mass $V$, missing
$\{R=-\infty\}$, set
\[
 I_\nu(u)=\int(R-u)d\nu,\qquad
 D_\nu=\sup_{u(v_0)=0}\{I_\nu(u)-B(u)\}.
 \tag{8.1}
\]
The difference is taken where $R$ is finite; no absolute integrability
of $R$ is stipulated. Finite $D_\nu$ implies integrability of every
genuine relative E1 test directly. Conversely if all are integrable,
their integrals are uniformly bounded on $B(u)\leq1$: otherwise choose
$I_\nu(u_j)\geq4^j$ and take the genuine decreasing convex sums
$\sum_{j\leq m}2^{-j}u_j+2^{-m}R$ of bounded cost. Their actual E1
endpoint has infinite relative integral by ordinary countable monotone
convergence. Contraction (6.7), with $r=(2^{n+1}B(u))^{-1/2}$ for
large cost, then bounds $I_\nu(u)$ by a constant times $1+\sqrt{B(u)}$,
so $D_\nu<\infty$. Thus (8.1) is an intrinsic relative energy condition.
It implies $\nu(Z)=0$ using the actual E1 pole constructed above.

Writing $C_0=2^{n+1}$, contraction also gives
\[
 I_\nu(u)\leq\epsilon B(u)+C_0D_\nu/\epsilon,
 \qquad 0<\epsilon\leq C_0.
 \tag{8.2}
\]
Hence near-maximizers have bounded actual energy cost. Section density
on $U$ and directed continuous Radon integration show that model-envelope
objectives have exactly the same supremum. The Radon net assertion used
here is elementary: increasing continuous superlevels on $U$ cover each
compact subset of the limiting open superlevel by finitely many sets;
directedness gives one dominating index, inner regularity gives its mass,
and finite layer-cake sums give the integral. No measurable-net monotone
convergence assertion is imported.

Choose model near-maximizers $a_j$ within $\varepsilon_j\to0$ and bounded
cost. Their outputs are finite divisorial measures. For a fixed continuous
test $h$, $z_{j,t}=P_\theta(a_j+th)$ lies between
$a_j\pm|t|\|h\|$. Its objective and obstacle inequality give
$E(z_{j,t})-E(a_j)-t\nu(h)\leq\varepsilon_j$.
Both signs of the uniform fixed-increment bound (7.3), first at fixed
$c,t$, then $j\to\infty$, then $t\downarrow0$, then $c\to\infty$,
prove $M(a_j)\to\nu$ weakly. This uses no attained inverse.

For fixed model envelopes $a,b,w$, (5.5) applied to every fixed majorant
of $(a+b)/2$ gives
$\limsup_m\delta_{A_m}(\psi_{m,a},\psi_{m,b})\leq\delta(a,b)$,
where $\delta(a,b)=E((a+b)/2)-(E(a)+E(b))/2$.
This is only an upper midpoint estimate; it does not assume the bounded
barycenter identification proved later. Formula (6.5), compact-uniform
bounded quotient convergence on $U$, and weak endpoint output convergence
then give the absolute model estimate. On a fixed continuous cutoff in
$U$, bounded changing tests pass by uniform convergence; omitted mass
is bounded by the limiting measure's omitted mass, since it misses $Z$.
Canonical output approximation and increasing continuous relative costs
extend it to all genuine tests and outputs:
\[
 \int|a-b|dM(w)\leq
 C_n\delta(a,b)^\alpha B(w)^{1/2}
       \max(B(a),B(b),B(w))^{1/2-\alpha}.
 \tag{8.3}
\]
First use $a=R$, bounded majorants of $b$, and the bound
$\delta(R,b_i)\leq B(b)/2$ to prove every actual relative test
integrable at $M(w)$. Then pass two test nets separately in absolute
$L^1$; their midpoint energies decrease. The same model-pair passage
through the pre-attainment outputs $M(a_j)\to\nu$ gives (8.3) at a
target, with the output cost bound of that sequence. These arguments
retain all generic terms in (6.5) before the degree limit.

Every output consequently satisfies the intrinsic target condition.
For fixed relative tests, (8.3) bounds the ordered majorant errors
uniformly over any bounded-cost output family. Fixed bounded relative
tests pass weak convergence on a fixed $U$ cutoff; then those absolute
errors pass the unbounded test. Contact is already forced by (4.7): at
$a=P_\theta(f_D)$ the equality $E(a)=\mathscr E(f_D)$ and $a\leq f_D$
give $0\leq\int(a-f_D)d\mu_D\leq0$. Thus finite planes, these exact
contact integrals, and the fixed-test passage give the whole supporting plane
\[
 E(v)\leq E(u)+\int(v-u)dM(u),\qquad u,v\in\EE.
 \tag{8.4}
\]
This is an outer-map plane, not the entire mixed-operator identification.

\section{Exact quotient traces and bounded physical identification}
Fix an actual $p:Y\to X$ with $p^*L=A+T$, $A$ nef and generated in
cleared degree, $T\geq0$. Multiplication by $s_{\ell T}$ injects the
\emph{full} $H^0(Y,\ell A)$ as a fixed subspace $S_\ell$ of
$R_\ell=H^0(X,\ell L)$. Let $V_A=A^n$, $\varepsilon=V-V_A\geq0$.
Use the induced subspace and quotient flags of the original (4.5) flag.
Extend an adapted subspace basis to a full adapted basis. Its trace is
exactly the sum of subspace and quotient traces. Uniform shifts by
$\ell\|f-g\|$ also shift every ordered quotient weight by at most this
amount. Quotient rank asymptotics therefore give
\[
 |\mathscr Q_A(f)-\mathscr Q_A(g)|\leq\varepsilon\|f-g\|,
 \qquad\mathscr Q_A=\mathscr E-\mathscr H_A.
 \tag{9.1}
\]
The same finite flag argument gives monotonicity. Rounding error is at
most the quotient rank and vanishes in the trace scaling.

Here is the required actual restricted-lattice calculation. Take a smooth
common equivariant total determination \emph{mapping to}
$Y\times\mathbb P^1$ and carrying the original obstacle $f_D$.
Write the total pulled $T$ as $T^{\rm hor}+T^{\rm ver}$. A generic
section in the restricted space has the literal form $s_{\ell T}a$,
$a\in H^0(Y,\ell A)$. In an actual local frame, its original extension
order is the order of $a$ plus the order of $s_{\ell T}$ in the pulled
frame. Along every central prime this is the extension condition in
\[
 p^*L+D-T^{\rm hor}=p^*A+D+T^{\rm ver}.
 \tag{9.2}
\]
Along horizontal primes, regularity as a section of the latter bundle
is precisely regularity of $a$ as a section of $A$ on its fixed generic
model. Thus both implications hold at every prime: multiplying gives
the original section with the prescribed horizontal zeros; division
gives an actual $A$ section and the retained central orders.
Normality identifies sections with these orders. Consequently this is
the full relative lattice of (9.2), not a selected series. This argument
requires the actual morphism to the specified auxiliary product; it
does not replace it by numerical equality. The example of blowing up
$(z,t)=(0,0)$ on $\mathbb P^1\times\mathbb P^1$ has
$p_X^*\{z=0\}=T^{\rm hor}+E$; the vertical $E$ term in (9.2)
is essential, and its cancellation confirms rather than contradicts
the formula.

If $V_A>0$, the time-twisted bundle (9.2) is big for sufficiently
large $c$. Its full lattice diagonalizes, and (4.2)--(4.4) applied to
this actual bundle give the restricted intercept $\mathscr H_A(f)$,
and its positive finite divisorial derivative $\mu_{A,T;f}$ of mass
$V_A$. Original and restricted derivatives test the same pulled original
model direction. The derivative of (9.1) has norm at most $\varepsilon$
and value $\varepsilon$ on the constant one, by the two constant laws.
A real functional with these properties is positive: for $0\leq h\leq1$,
the norm bound on $1-h$ gives its value on $h$ at least zero.
Model density extends this to continuous tests. Hence
\[
 0\leq\mu_f-p_*\mu_{A,T;f},\qquad
 \|\mu_f-p_*\mu_{A,T;f}\|_{\rm TV}=V-V_A.
 \tag{9.3}
\]
If $V_A=0$, restricted rank is $o(\ell^n)$ and its fixed-obstacle
weights are bounded linearly in $\ell$, so its trace and derivative
are zero directly. No differential at a nonbig class is used. In the
moving setting a cofinal tail has $V_m>0$ by (5.4).

Suppose a bounded rational generated scalar $-C\leq\xi_m\leq0$
on $A_m$ has actual physical FS stage $v_m=q_{T_m}+\xi_m$, and
$f\geq v_m$ is an original model obstacle with $f(v_0)\leq0$.
All full restricted weights lie in $[-\ell C,0]$: the lower bound
uses $f\geq v_m$ and the canonical norms of auxiliary sections, and
the upper bound evaluates at $v_0$. Taking powers of the rational
weighted generators of $\xi_m$, after clearing their weights, shows
that the complete restricted stages $\sigma_{\ell,f}$ satisfy
$\xi_m\leq\sigma_{\ell,f}\leq0$ and
$\sigma_{\ell,f}\leq f-q_{T_m}$ at function-field divisorial points.
The fixed twist $c>C+1$ makes each compactified complete system
globally generated. Its full section bundle is the fixed big class
$\beta=p^*A_m+D+T_m^{\rm ver}+cF$.
Complete-base-system Fujita gives its nef moving top volumes tending
to $\vol(\beta)$. Subtracting the fiber term proves
\[
 e_m(\sigma_{\ell,f})\to\mathscr H_{A_m}(f).
 \tag{9.4}
\]

Volume convergence identifies every fixed derivative test as well.
Let $\Gamma_\ell\leq\beta$ be these complete nef models. They satisfy
$\Gamma_\ell\leq p^*L+cF$. For a fixed original model direction $H$,
choose a fixed globally nef big $\omega$ above this last class with
$\omega\pm H$ nef, on a fixed determination carrying $H$, and pull
it to the common refinements. The numerical upper bound suffices for
BFJ Corollary 3.4; the actual effective difference $\beta-\Gamma_\ell$
also gives volume order. For both signs of $t$ the independently read
Morse bound gives
\[
 \vol(\beta+tH)\geq\Gamma_\ell^{n+1}
             +(n+1)t\Gamma_\ell^nH-Kt^2.
 \tag{9.5}
\]
The same $K$ works for all $\ell$. Pass top volumes first; the positive
and negative signs bound limsup and liminf of $\Gamma_\ell^nH$.
Then let $t\to0$. BFJ differentiability gives
$\Gamma_\ell^nH\to\langle\beta^n\rangle H$.
For vertical tests this is exactly the finite restricted derivative.
Model density and its fixed raw mass therefore identify the weak limit
of its finite top measures with $\mu_{A_m,T_m;f}$. This proves the
complete two-sided derivative squeeze, without an unbounded auxiliary
operator or uniform Fujita rate across changing obstacles.

Finite concavity at $\xi_m$ and (9.4) give, with
$\sigma_m=\MA_{A_m,\rm raw}(\xi_m)$,
\[
 0\leq\mathscr H_{A_m}(f)-e_m(\xi_m)
 \leq\int\min(C,f-v_m)d\sigma_m.
 \tag{9.6}
\]
For a fixed original test $h$, ordered finite energy telescoping gives
$I\leq(n+1)$ times the ordered gap. A fixed-$C$ twist makes all
bounded auxiliary bundles nef and dominated by $p^*L+cF$.
Writing $H$ as two fixed nef tests bounds its quadratic forms in these
slots, independently of $m,f$ and refinements. Cauchy--Schwarz gives
\[
 |p_*\mu_{A_m,T_m;f}(h)-p_*\sigma_m(h)|
 \leq C_{h,C}\sqrt{\mathscr H_{A_m}(f)-e_m(\xi_m)}.
 \tag{9.7}
\]
The constant is fixed before the moving limit.
The same fixed twist bounds $\int v(S)d\sigma_m$ uniformly for
every fixed effective original divisor $S$: its vertical part is the
finite moment, its horizontal part is effective, and a fixed nef test
bounds the total intersection by a positive product of $p^*L+cF$.
Thus auxiliary collar masses are at most $C_{S,c}/H$.
Original bounded-cost outputs have the separate collar bound
$C_{S,c}/H+D/c$ from (7.6). Both collars must be kept distinct.

\begin{theorem}[Actual bounded generated quotient identification]\label{bounded}
Let $w\in\EE$ be normalized, continuous on $U$ and
$R-C\leq w\leq R$. Suppose $-C\leq\xi_m\leq0$ are rational
generated quotients on the common complete-system $A_m$, their generic
values tend to zero, their actual original FS stages $v_m$ are at most
$w$ on all lifts, and $\xi_m\to w-R$ uniformly on each compact subset
of $U$. Then $e_m(\xi_m)\to E(w)$ and $p_*\sigma_m\to M(w)$ weakly.
\end{theorem}
\begin{proof}
Choose fixed integer $N>C+2$ and genuine floors
$u_m=\max(v_m,R-N)\leq w$. They are actual E1 functions of bounded
normalized cost at most $NV$; their generic values $b_m\to0$ are
retained. On $U$, $0\leq w-u_m\leq N$, with uniform convergence
to zero on compact subsets. Original collars and the whole plane (8.4)
give $E(u_m)\leq E(w)\leq E(u_m)+o(1)$. Ordered energy variation
(7.5) gives $M(u_m)\to M(w)$.

At fixed $m$ take nonpositive canonical model majorants $f_j\downarrow
u_m$ with generic value $b_m$. Every original quotient weight is
nonpositive, so $\mathscr E(f_j)\leq\mathscr H_{A_m}(f_j)$.
Pass $j$ on the fixed finite support of $\sigma_m$ in (9.6):
\[
 E(u_m)\leq e_m(\xi_m)
       +C\sigma_m\{v_m<R-N\}.
 \tag{9.8}
\]
Auxiliary collars make this last mass tend to zero: on any fixed compact
subset of $U$, $v_m\to w\geq R-C>R-N$ with a fixed margin.
First fix the collar, then the moving index. The reverse bound is (5.5),
holding every finite original majorant of $w$ fixed through the degree
and moving limits before its outer infimum. Thus energies converge.

Combine (9.3), (9.7) and the same fixed-support limit to get
\[
 |M(u_m)(h)-p_*\sigma_m(h)|
 \leq(V-V_m)\|h\|+C_{h,C}\sqrt{C\sigma_m\{v_m<R-N\}}.
 \tag{9.9}
\]
Both terms vanish. The weak floor output limit identifies every model
test; density and raw mass identify all continuous tests. No original
E1 membership is assigned to the physical $v_m$.
\end{proof}

This theorem applies to maxima and rational convex combinations of any
fixed finite list of model envelopes. Tensor their actual generators
after clearing rational coefficients. Their physical horizontal factor
is exactly $T_m$ because the coefficients sum to one. Their compact
quotient limits and uniform boundedness follow from (5.3)--(5.4).
In particular it proves the exact bounded midpoint energy identification
needed for the atomic density argument, rather than only the upper estimate preceding (8.3).

\section{The genuine mixed operator and its entire energy identity}
Choose once a finite rational unisolvent set on the simplex for homogeneous
degree $n$ measure polynomials and degree $n+1$ energy polynomials.
It exists because a nonzero polynomial cannot vanish on every rational
interior point; choose evaluations successively spanning the finite
dual space. At every such fixed weight Theorem~\ref{bounded} passes
the finite nef polynomial to the original diagonal $M$ and energy $E$.
Invert the fixed matrices before any moving limit. Every finite mixed
coefficient is a positive measure of raw mass $V_m$, retaining the
coefficient $b_E$ times its restriction intersections. The limits are
positive symmetric measures of mass $V$ and independent of the complete
system choice. Every fixed real simplex weight follows as well: bounded
relative order gives uniform parameter continuity of energy and its
perturbed secants, hence weak continuity of the diagonal output.

For arbitrary genuine endpoints choose independent canonical model
envelopes decreasing to them. At each fixed rational evaluation weight
their combinations decrease genuinely; (7.4) and the secant brackets at
$-t,0,t$ prove decreasing weak continuity of their diagonal outputs.
The same fixed inverse matrices give positive symmetric original mixed
measures, denoted $\MA_\theta(u_1,\ldots,u_n)$, with diagonal $M$.
Their expression using only diagonal outputs proves independence of
canonical choices and decreasing weak continuity along every genuine
finite-energy net, not just one canonical approximation. All fixed real
weights pass at the majorant limit without exchanging parameter and net.
Positive finite coefficients give
\[
 0\leq\MA_\theta(u_1,\ldots,u_n)
 \leq \frac{n^n}{n!}M\left(\frac{u_1+\cdots+u_n}{n}\right).
 \tag{10.1}
\]
Consequently every mixed measure integrates every actual relative E1
test and misses $Z$. This is the genuine operator: it is constructed
from finite actual intersection products, their bounded full-system
limits and genuine decreasing endpoints on the original carrier. It is
not assigned by an inverse and not an abstract completion.

For mixed decreasing endpoints with finite-energy limits, their average
costs are bounded on a tail. Domination (10.1) and (8.3) make the ordered
errors of any canonical test majorant uniformly small at all these mixed
outputs. At a fixed majorant its relative cost is bounded continuous on
$U$; hold an original continuous cutoff fixed in the weak limit, then
remove it using Radon mass. Finally remove the majorant by those uniform
absolute errors. This proves weighted convergence for every fixed genuine
relative test and continuous multiplier. For an independently decreasing
genuine test use its ordered midpoint bound
$\delta(v_k,v)\leq(E(v_k)-E(v))/2$ and generic normalization error
$V|v_k(v_0)-v(v_0)|\to0$. This proves the full product-directed weighted
limit, including changing tests, without measurable-net MCT.

For two bounded model envelopes the finite nef telescoping identity passes
through Theorem~\ref{bounded} and the coefficient inversion. Its test
difference is bounded, converges compact-uniformly on $U$, and the mixed
auxiliary collar follows from the same fixed nef intersection bound as
in Section 9. Hold that collar first. For arbitrary genuine endpoints
take their independent model-envelope nets and use the just-proved weighted
continuity and the sum of their two uniform ordered errors. Thus
\[
 E(u)-E(v)=\frac1{n+1}\sum_{j=0}^n
 \int(u-v)d\MA_\theta(u^{[j]},v^{[n-j]}).
 \tag{10.2}
\]
All are finite original relative integrals, with constants restored by
the raw mass $V$. This constructs the entire actual E1 mixed operator
and energy identity. Together with Sections 2--8 it discharges S1--S3
and M1--M2 rather than assuming them as an axiomatic package.

\section{Reconstruction of the finite unbounded truncation route}
The bounded physical identification and genuine mixed map used here have
been proved above. The finite estimates are reconstructed next, retaining
generic errors before every limit.

For a nonpositive finite generated scalar $\xi$ on a nef $A$ and
$\xi^C=\max(\xi,-C)$, strict finite nef locality follows from G0's
finite ample locality by adding the same actual $\epsilon H$ zero metric
and passing its finite component polynomials. The finite energy identity
has nonnegative terms and gives
$\int(-\xi)d\mu_\xi\leq(n+1)B_A(\xi)$.
Locality identifies $\mu_\xi$ and $\mu_{\xi^C}$ on $\{\xi>-C\}$.
Equal raw masses imply equal complementary masses, so Markov gives
\[
 \|\mu_\xi-\mu_{\xi^C}\|_{\rm TV}
       \leq 2(n+1)B_A(\xi)/C.
 \tag{11.1}
\]
For $\eta\leq\xi\leq0$ of costs at most $D$, the plane at the
lower cap, its cost monotonicity and (6.5) give
\[
 0\leq e_A(\xi^C)-e_A(\eta^C)
 \leq V_A|g_\xi-g_\eta|
       +K_nD^{1-\alpha}(e_A(\xi)-e_A(\eta))^\alpha.
 \tag{11.2}
\]
No quadratic estimate involving an unbounded auxiliary slot is used.
At each fixed cap a twist $c>C+1$ globally generates the compactified
bounded system: discard weights below $-mC$, include the weight-zero
generated generic system at floor $-mC$, and supply all time monomials.
Its nef bundle is effectively at most $p^*L+cF$.
Write a fixed model test $H$ as two fixed globally nef classes.
Positive-product order bounds its quadratic form in any bounded slots
uniformly. Finite ordered $I\leq(n+1)$ times the energy gap and
Cauchy--Schwarz give a fixed-$C$ continuous-test bound by the square
root of that gap.

An energy-complete joint filter imposes only: obstacle indices are
cofinal; degrees are nonempty, cleared and divisibility-cofinal;
individual auxiliary errors
$e_m(\psi_{m,f_i})-E(P(f_i))$ and generic errors tend to zero.
It imposes no output convergence. It exists from (5.4): for each
finite tuple, tolerance and prior degree, select a sufficiently large
common multiple with every finite-list error below the tolerance.
The valid tuples, ordered by dominating indices, divisibility and
decreasing tolerance, are directed. For a fixed canonical upper index
$k$, full filtration order gives the moving lower scalar below
$\psi_{m,f_k}$. At fixed cap (11.2) and the fixed-$C$ quadratic bound
compare their measures. First take the moving degree limit, retaining
the generic error until it vanishes; then take $k$ to its endpoint;
finally remove the cap by (11.1). Theorem~\ref{bounded} identifies the fixed upper cap with
$M(\max(P(f_k),R-C))$. Its genuine decreasing limit as $k$ tends to
the endpoint is $M(\max(u,R-C))$. At fixed $C$, (11.2) and the uniform
quadratic bound make the lower cap have that same weak limit. At every
fixed original model test the limiting error is at most a constant
times $(E(P(f_k))-E(u))^{\alpha/2}$, after the generic term vanishes.
Then (11.1) bounds the untruncated error by
$2(n+1)D\|h\|/C$. Finally $\max(u,R-C)\downarrow u$ genuinely;
decreasing weak continuity removes $C$. This proves every energy-complete
diagonal weak limit independently of degree choices.

For rational $t_a\geq0$, $\sum_a t_a=1$, write
$\xi_{\lambda,t}=\sum_a t_a\xi_{\lambda,a}$ and
$\xi_{m,k,t}=\sum_a t_a\psi_{m,f_{a,k_a}}$. After clearing powers
both are generated systems on the same $A_m$, and their physical
horizontal factor is precisely $T_m$. The finite plane at the lower
barycenter and (6.5) give
\[
\begin{aligned}
0&\leq e_m(\xi_{m,k,t})-e_m(\xi_{\lambda,t})\\
&\leq\sum_a t_a\bigl[ V_m|g_{m,k_a}-g_{\lambda,a}|
 +K_nD^{1-\alpha}
 (e_m(\psi_{m,f_{a,k_a}})-e_m(\xi_{\lambda,a}))^\alpha\bigr].
\end{aligned}
\]
The lower barycenter cost is at most $D$ by concavity. At fixed $k$
Theorem~\ref{bounded} identifies the upper energy; first pass the
moving index, then $k$ to its endpoint. The right side tends to zero
and the upper genuine barycenters decrease with their energies. Thus
individual energy completeness implies barycenter energy completeness;
it is not an extra admissibility condition. A fixed finite rational
unisolvent matrix for the degree-$n$ raw measure polynomial and
degree-$(n+1)$ energy polynomial then identifies mixed coefficients.
This uses only finite nef multilinearity. It never interpolates ordinary
big positive products as multilinear objects.

For an independently changing physical test quotient $\eta_\lambda$,
finite positive mixed domination by $n^n/n!$ times its average diagonal
output and (6.5) compare it with one fixed canonical upper test.
More precisely, after the generic errors vanish at fixed upper index $k$,
\[
\limsup_\lambda\int(\eta_{m,k}-\eta_\lambda)d\sigma_\lambda
\leq\frac{n^n}{n!}K_nD^{1-\alpha}
 (E(P(g_k))-E(v))^\alpha .
\]
At that fixed upper index the quotient is bounded; a fixed original
$U$ cutoff passes its changing test through the weak output limit.
Then the ordered genuine-test errors tend to zero. The resulting
conclusion is the paired moment
\[
 \int(h\circ p_m)(-\eta_\lambda)d\sigma_\lambda
 \longrightarrow\int h(R-v)d\MA(u_1,\ldots,u_n).
 \tag{11.3}
\]
At finite divisorial atoms it is exactly the difference between the
physical reference $q_{T_m}$ and physical paired test
$q_{T_m}+\eta_\lambda$. It is not the fixed original test $R-v$.
Writing $v_{\rm phys}=q_{T_m}+\eta_\lambda$, their difference is
\[
 (R-v)-(-\eta_\lambda)=(R-q_{T_m})+(v_{\rm phys}-v).
 \tag{11.4}
\]
Neither term is automatically small in absolute physical-output
integral. At supported poles (11.4) is not used as infinite arithmetic.
We do not claim fixed original unbounded weighted convergence at every
physical auxiliary output. Genuine canonical weighted continuity was proved in Section 10 and
the necessary energy-complete weak approximation independence was proved
in this section. Neither proof uses the stronger physical assertion.
No uniform absolute control forcing the defect (11.4) to zero has been
proved here from the stated admissibility. The stronger assertion remains
unproved; this is not a counterexample to that assertion.

\section{Atomic dual density before attainment and the actual equation}
The finite four-point estimate on the same actual nef $A$ is
\[
\begin{aligned}
 I_A(\xi,\eta)&=\sum_{j=0}^{n-1}
 \|\xi-\eta\|^2_{T_\xi^jT_\eta^{n-1-j}}\geq0,\\
 \left|\int\zeta(d\mu_\xi-d\mu_\eta)\right|
 &\leq C_n I_A(\xi,\eta)^{1/2}J_A(\zeta)^\alpha
       \max(J_A(\zeta),J_A(\xi),J_A(\eta))^{1/2-\alpha}.
\end{aligned}
 \tag{12.1}
\]
Top-power telescoping gives the first identity. Cauchy--Schwarz in its
$n$ graph forms and (6.4) for the test relative to zero give the second.
Constants cancel because the two raw output masses agree; there is no
generic term in this estimate. Equations (6.3) and $I_A\leq n d_A$ give
$I_A\leq4^{n+1}\delta_A$.

For fixed model envelopes use the same full-system $A_m$ and horizontal
factor. Their bounded quotients pass all test differences on a fixed $U$
cutoff, with omitted mass controlled by the weak output limit which misses
$Z$. Their $J_A$ values tend to genuine costs, retaining generic errors.
Theorem~\ref{bounded} identifies their midpoint energy exactly.
Then independent genuine canonical weighted limits and decreasing energies
pass the assertions to all actual normalized $a,b,z\in\EE$. Thus with
$I(a,b)=\int(a-b)(dM(b)-dM(a))$,
\[
\begin{aligned}
 \left|\int(R-z)(dM(a)-dM(b))\right|
 &\leq C_n I(a,b)^{1/2}B(z)^\alpha
     \max(B(z),B(a),B(b))^{1/2-\alpha},\\
 0\leq I(a,b)&\leq4^{n+1}\delta(a,b).
\end{aligned}
 \tag{12.2}
\]
These are actual relative integrals, finite by (8.3); no target inverse
has been used. For $n=1$ the last exponent-zero factor is omitted.

Let $\nu$ be any intrinsic target (8.1). Section 8 constructed normalized
model near-maximizers $a_j$, costs at most $D_0$, objective gaps
$\varepsilon_j\to0$, and finite divisorial outputs $\nu_j=M(a_j)$
already weakly converging to $\nu$. Their midpoint objective is at most
the intrinsic supremum, so
$\delta(a_j,a_k)\leq(\varepsilon_j+\varepsilon_k)/2$.
For every finite $D$, (12.2) implies uniformly on actual normalized
$B(z)\leq D$,
\[
 |\int(R-z)(d\nu_j-d\nu_k)|
 \leq L_{n,D,D_0}\sqrt{\varepsilon_j+\varepsilon_k}.
 \tag{12.3}
\]
Identify its limit at each fixed $z$ without passing an unbounded weak
test. First pass a fixed bounded canonical test majorant on a fixed
original $U$ cutoff through $\nu_k\to\nu$. Ordered errors from (8.3)
are uniform over the bounded-cost outputs $\nu_k$ and tend to zero.
At the fixed intrinsic target the increasing continuous relative costs
converge by the directed Radon finite-cover argument of Section 8.
Hence $\int(R-z)d\nu_k\to\int(R-z)d\nu$ for each fixed $z$.
Let $k\to\infty$ in (12.3) separately for each $z$; its constant is
independent of $z$, so only afterwards take the supremum. This proves
\[
 \sup_{B(z)\leq D}|\int(R-z)(d\nu_j-d\nu)|
       \leq L_{n,D,D_0}\sqrt{\varepsilon_j}\longrightarrow0.
 \tag{12.4}
\]
The finite atomic density precedes attainment and assumes no inverse.

Every $z\mapsto\int(R-z)d\nu_j$ is continuous in divisorial coordinates:
its finite support has finite genuine values by the actual ample embedding.
Its uniform convergence (12.4) on each normalized energy sublevel proves
net continuity of the intrinsic target integral there. Those sublevels
are closed in the genuine normalized compact cone by energy weak usc.
Coercivity puts a model maximizing sequence in one of them. A convergent
subnet has an \emph{actual} original finite-energy limit $u$.
Target-integral continuity and energy weak usc make its objective at least
the limiting supremum, hence exactly that supremum. This is actual weak
attainment; neither an energy completion nor a lift supplies the solution.

For $h\in C(K)$ and real $t$, $u_t=P_\theta(u+th)$ is actual E1 and
$u_t\leq u+th$. Its objective is constant invariant because both raw
constant factors are $V$. Maximality gives
\[
 0\geq F_\nu(u_t)-F_\nu(u)
 \geq E(u_t)-E(u)-t\int h\,d\nu.
 \tag{12.5}
\]
Both signs of (7.4) give $M(u)(h)=\nu(h)$ for every original continuous
test. Riesz uniqueness gives $\MA_\theta(u)=M(u)=\nu$ on $K$.
The target has not defined the operator.

Every solution, not just this maximizer, is unique modulo constants.
Whole plane (8.4) makes every solution for $\nu$ maximize $F_\nu$.
Two normalized maximizers $a,b$ have $\delta(a,b)=0$ by midpoint
maximality and integral linearity. For every rational divisorial $z$,
$V\delta_z$ is an intrinsic target independently of its inverse:
the uniform coordinate bound of Lemma~\ref{embed} bounds every relative
integral and its dual supremum. The target version of (8.3), proved
using pre-attainment weak output approximation, then gives $a(z)=b(z)$.
All-point ample order determination gives $a=b$ everywhere, including
supported poles. Removing generic normalization restores one constant.
Conversely (8.3) puts every operator output in the intrinsic target class,
and its whole plane computes the dual supremum. Thus (8.1) is the exact
target condition, not a chosen sufficient support condition.

\section{Actual strict Borel locality and comparison}
For model envelopes $a=P(f)$, $b=P(g)$, the auxiliary maximum of their
bounded quotients satisfies Theorem~\ref{bounded}, with endpoint
$w=\max(a,b)$. It is an actual FS stage after multiplying the same
horizontal factor, below $w$ at every lift, and compact-uniformly tends
to $w-R$. Thus its raw measure tends to $M(w)$.
On a compact subset of $\{a>b\}\cap U$ there is a positive scalar gap.
Finite nef locality therefore identifies its measures with those of $a$
on every continuous test supported there for a degree tail. Pass weak
limits and apply Radon uniqueness on that open subset of $U$.
This proves bounded strict locality, with the symmetric identity for $b$.
There is no assumed identity $\max(P(f),P(g))=P(\max(f,g))$.

The bounded identity also allows independently shifted input constants.
Subtract their maximum generic value from both inputs. The maximum is
then normalized, whereas the two inputs may have negative generic values.
Their quotient systems are shifted by exactly the corresponding constants;
they remain bounded above by zero, and their maximum has generic limit zero.
Its physical stage still satisfies Theorem~\ref{bounded}. The endpoint
measures are unchanged by the constant law, and finite locality has the
unchanged scalar winner. This proves the bounded identity in that case too.

For arbitrary $u,v\in\EE$ make this same common shift and take independent canonical model
envelopes $u_i\downarrow u$, $v_i\downarrow v$, and
$w_i=\max(u_i,v_i)\downarrow w=\max(u,v)$. Costs of all outputs
are bounded by one $D$. Their genuine decreasing weak continuity follows
from the fixed energy secant brackets at $-t,0,t$ in (7.4).
Canonical majorants keep each input's generic value exactly fixed;
ordered test errors therefore have zero generic difference. Costs in the
estimates mean the costs after individual generic normalization, and the
midpoint defect is invariant under those individual constants.
Equation (8.3) gives uniform ordered test errors
\[
 \sup_{B(z)\leq D}\int(u_j-u)dM(z)=:\eta_j(D)\to0,
 \tag{13.1}
\]
with an analogous bound for $v$. The bound uses
$\delta(u_j,u)\leq(E(u_j)-E(u))/2$ and is independent of the output.

Let $\beta$ be zero on $(-\infty,1]$, one on $[2,\infty)$ and linear
between. Fix $h\in C(K)$ with compact support in $U$, then fix $j$ and
$\epsilon>0$. The fixed continuous test on $K$ is
$h\beta((u_j-v_j)/\epsilon)$, extended by zero off $U$.
For $i\geq j$, if its ramp is positive but $u_i\leq v_i$, then
$u_j-v_j>\epsilon$, $v_i\leq v_j$, hence $u_j-u_i>\epsilon$.
On $\{u_i>v_i\}$ bounded locality cancels the two measures. On its
complement, (13.1) bounds both positive offending measures and gives
\[
 |\int h\beta((u_j-v_j)/\epsilon)d(M(w_i)-M(u_i))|
       \leq2\|h\|\eta_j(D)/\epsilon.
 \tag{13.2}
\]
Pass the weak limit at fixed $j,\epsilon,h$. At the now fixed limiting
measures, both relative costs are finite almost everywhere by (8.3).
The Lipschitz ramp error is bounded by
$\epsilon^{-1}((u_j-u)+(v_j-v))$. Its absolute integrals tend to zero
by (13.1). Let $j$ tend to its endpoint in (13.2), and only then take
an ordinary sequence $\epsilon\downarrow0$. Dominated convergence at
the two fixed measures selects $\mathbf1_{\{u>v\}}$.
Radon uniqueness and nullness of $Z$ give the original Borel identities
\[
 \mathbf1_{\{u>v\}}M(\max(u,v))=\mathbf1_{\{u>v\}}M(u),
 \quad
 \mathbf1_{\{v>u\}}M(\max(u,v))=\mathbf1_{\{v>u\}}M(v).
 \tag{13.3}
\]
No unbounded winner is assumed open and no changing indicator is used
as a weak test. The diagonal identity is the locality assertion needed
for the equation and comparison below.

Apply (13.3) to $\max(u,v-\epsilon)$. Raw masses are $V$ and constants
do not change the operator. Consequently
$M(v)\{u<v-\epsilon\}\leq M(u)\{u\leq v-\epsilon\}$.
Both sets increase to $\{u<v\}$ outside the null infinite-relative-cost
sets as an ordinary sequence $\epsilon\downarrow0$. Finite measure
continuity from below proves
\[
 \int_{\{u<v\}}d\MA_\theta(v)\leq
       \int_{\{u<v\}}d\MA_\theta(u).
 \tag{13.4}
\]

\section{Negligible sets and the additional zero-volume theorem}
There is an unconditional geometric pole-covering deduction in the
retained big domain. Embed every member of a bounded-above raw-supremum
family into a common actual ample $\Omega$ as the identical scalar.
Its raw supremum and its usc regularization are literally identical on
the same topological carrier, including supported points. Thus its
discrepancy is ample-negligible in the primary definition. BJ 13.17,
whose full pertinent proof was independently read with its ample EP
hypotheses, supplies ample-psh $p$ with a pole covering that discrepancy.
Choose rational $t>0$ with $L-t\Omega$ big and a section giving the
actual identity $L=t\Omega+E$. Tensor G0's fixed-bundle decreasing FS
net for $p$ with $s_E$. The one actual net $tp_i+q_E$ has a genuine
big-class endpoint covering the same discrepancy. No equality merely
on divisorial coordinates, arbitrary lift or subtraction of poles is used.
Section 6 makes an actual E1 pole covering it outside the reference poles.
Every intrinsic target, and every output integrating all genuine relative
tests and missing reference poles, consequently has a Borel null superset
of each such discrepancy. This proves outer nullness even if the original
negligible subset is not assumed Borel.

\begin{theorem}[Entire numerical-zero range]\label{zero}
Let the actual rational $L$ be numerically trivial, including nontorsion
$\operatorname{Pic}^0$ bundles, on the retained $X$. For the same varying
actual-bundle single-net cone at zero and zero scalar reference,
$\PSH(0)=\mathbb R$. Its entire finite-energy space is $\mathbb R$,
its raw energy and raw top operator are zero, and its sole positive
mass-zero target is zero. The equation is attained and unique modulo
constants; comparison and both derivative signs are trivial.
\end{theorem}
\begin{proof}
Every actual $L+\epsilon A$ is ample. Its zero-weight generating scalar
is exactly zero everywhere. These constant scalars with classes tending
to zero give one defining net; shifts give all constants.
Normalize $u\in\PSH(0)$ at $v_0$. Integer tensoring gives $mu\in\PSH(0)$
for every $m\geq1$. Embed all of them into one fixed actual ample cone.
Its uniform coordinate bound gives $-C_z\leq mu(z)\leq0$ for every
$m$, hence $u(z)=0$. G0's all-point determination gives $u=0$.
The constant ample perturbations have raw masses $\epsilon^n(A^n)\to0$.
Positivity bounds every continuous test by that mass times its sup norm,
so the unique raw measure limit is zero. Their constant raw energies
$c\epsilon^n(A^n)$ tend to zero. This defines the zero energy on the
entire classified cone. A positive measure of zero mass is literally
zero, and the quotient of constants by constants has one element.
No zero mass is probability-normalized.
\end{proof}

This does not prove the nonzero zero-volume pseudoeffective range.
On $\mathbb P^1\times\mathbb P^1$, $L=p_1^*\OO(1)$ has volume zero,
and the actual generated-system scalar with coordinate weights $0,-1$
is nonconstant. It and zero have raw top perturbation masses
$2\epsilon(1+\epsilon)\to0$, so both top limits vanish. Hence uniqueness
modulo constants fails for that whole bounded cone with the raw zero
top equation. A different replacement datum is required for any proposed
larger zero-volume theory.

\section{Canonical result, retained domain and remaining stronger question}
\begin{theorem}[Retained original P16.5, entire first big domain]
Under the field and geometric hypotheses of Section 1, for every actual
big rational $L$, the cone of Section 2 has genuine original decreasing
closure, envelopes and normalized divisorial compactness. Its intrinsic
reference is $R=P_\theta(0)$, and its entire finite-energy space is (4.6).
There are positive symmetric genuine original mixed Radon measures on this
space, of raw mass $V=\vol_X(L)$, with model coefficients, all genuine
decreasing weak and weighted continuity, and energy identity (10.2).
They have the energy-complete full-system weak and paired relative weighted
independence of Section 11. The original continuous outer-energy derivative
is (7.4). Strict Borel locality and comparison are (13.3)--(13.4).

A positive original Radon measure $\nu$ is a target precisely when it has
mass $V$, misses $\{R=-\infty\}$ and has finite intrinsic dual energy
(8.1), equivalently integrates every actual relative E1 test. No absolute
integrability of $R$ or additional support condition is imposed. Every such
target has an actual original finite-energy solution, and every solution
is unique modulo one real constant. The additional zero-volume domain
proved is exactly the entire actual numerically trivial rational-bundle
range of Theorem~\ref{zero}, with zero energy and literal zero target.
\end{theorem}
\begin{proof}
Sections 2--3 prove the genuine scalar closure, envelopes and compactness.
Sections 4--7 construct the full outer energy, finite products and original
derivative map with both signed big margins. Section 5 proves every finite
list full-system limit with actual horizontal factors and raw multiplicities;
Section 6 proves contraction and an actual E1 divisor pole. Section 8 proves
all genuine relative-test integrability and whole planes before an inverse.
Sections 9--10 prove the exact quotient trace and two-sided restricted
derivative squeeze, bounded generated quotient identification, entire genuine
mixed map and energy identity. Section 11 proves the necessary unbounded
approximation independence. Section 12 proves atomic dual density before
actual weak attainment, the attained equation and every-solution separation.
Section 13 proves genuine Borel comparison by fixed ramps. The zero-volume
classification and equation are proved separately in Theorem~\ref{zero}.
Thus no conditional S--M package remains in this retained first domain.
\end{proof}

For $n=0$, the retained integral projective carrier is one point; every
bundle is ample, the mass is one, and G0 gives the trivial equation. This
does not hide a positive-dimensional nonample case.

The stronger question still open in this draft is whether every fixed
original unbounded relative E1 test has weighted convergence at every
physical auxiliary output under the energy-complete conditions alone.
It requires controlling the two separate terms of (11.4) in absolute
physical-output integral. It is not asserted by paired quotient moments,
and no infinite difference is evaluated at supported poles. It is also
not used to define the genuine operator, prove its genuine canonical
weighted continuity, establish the necessary weak approximation
independence, or attain and separate the original equation. No further
nonzero zero-volume pseudoeffective range is proved. The positive theorem
is a full-domain unverified mathematical draft in the retained first big
domain, with explicit additional zero range; its canonical result is
\texttt{candidate=true} in that exact assignment scope.

\appendix
\section{Evidence and reproducibility}
The exact supplied JSON bytes are retained in \texttt{frozen/input-original.json},
with decoded task and assignment preserved separately. All sixteen marked
whole bodies are retained in \texttt{supplied/}, with literal inclusive
assignment-LF coordinates and fresh byte hashes. Old ledger paths were
not opened. G0/G1 are the only mathematical grants; unreviewed papers
are not converted into premises by their historical descriptions.

Independent public reads use BJ arXiv:1801.08229v4 (24 March 2022),
BFJ math/0608260v2 (4 September 2007), and Addendum 2206.07183v2
(7 September 2023). Complete units retained with no missing renderer
line are: BFJ310--441 and445--545; BJ2705--3107,3996--4022 and
10614--10647; Addendum128--257. Other ancillary primary readings
are retained in the raw responses. The source register distinguishes
applied algebraic premises, the primary ample separator and negligible
theorem, source matching within G1, and new deductions. No whole-primary
or recursive proof reading is claimed. BFJ's automatic HTML date does
not change its version. Literal renderer labels and PDF page labels are
preserved separately from inclusive one-based LF lines of extracted text.
The attempted native PDF download failed at the inherited proxy; no
downloaded original PDF or its hash is invented.

Run \texttt{python3 rerun.py} in this workspace. It verifies frozen inputs,
all whole-unit slices, primary maps, hashes, result schema and finite
rational checks, and compiles this standalone paper twice with the installed
pdfLaTeX. Commands, outputs, versions and artifact hashes are retained.
Finite checks concern transcription and explicit examples; compilation
and execution do not establish mathematical correctness.
No subagents, other campaigns, private outputs, credentials, authentication
or runtime settings were used or changed. No resource-heavy computation
was launched; no operator-provided supervised compute service was exposed.

\begin{thebibliography}{9}
\bibitem{BJ} S. Boucksom and M. Jonsson, \emph{Global pluripotential theory
over a trivially valued field}, \href{https://arxiv.org/pdf/1801.08229v4}{arXiv:1801.08229v4}.
\bibitem{BFJ} S. Boucksom, C. Favre and M. Jonsson, \emph{Differentiability
of volumes of divisors and a problem of Teissier},
\href{https://arxiv.org/html/math/0608260v2}{arXiv:math/0608260v2}.
\bibitem{ADD} S. Boucksom and M. Jonsson, \emph{Addendum to global
pluripotential theory over a trivially valued field},
\href{https://arxiv.org/pdf/2206.07183v2}{arXiv:2206.07183v2}.
\end{thebibliography}
\end{document}

% editorial-numbering-final
